% Chapter 5 — BiDi: Hebrew and English in Multi-Agent Systems
\addtocontents{toc}{%
  \protect\etocsetstyle{chapter}{}{}{\protect\addvspace{1.0em plus 1pt}%
   \begingroup\parindent0pt\noindent\bfseries%
   \etocpage\hfill{\color{uohblue}\etocname\hspace{1em}\etocnumber}%
  \par\endgroup}{}%
  \protect\etocsetstyle{section}{}{}{\begingroup\parindent0pt\noindent\normalfont\normalsize%
     \etocpage\hspace{0.5em}\protect\hebrewdotleaders\hspace{0.5em}%
     \etocname\hspace{0.5em}\etocnumber%
  \par\endgroup}{}%
}
\hebrewheadingformat
\chapter{\texthebrew{דו-כיוניות: עברית ואנגלית במערכות סוכנים}}
\label{ch:bidi}

% ---------------------------------------------------------------
\section[\texthebrew{משמעות הדו-כיוניות}]{\texthebrew{משמעות הדו-כיוניות}\mbox{\setLTR\enspace\thesection}}
\label{sec:bidi-meaning}

\begin{hebrew}
עיבוד טקסט דו-כיווני הוא אחד האתגרים המרכזיים בפיתוח מערכות תוכנה רב-לשוניות.
אלגוריתם הדו-כיוניות של תקן \foreignlanguage{english}{Unicode} קובע כיצד מערכות
תצוגה אמורות לעבד ולהציג טקסט המכיל שפות הנכתבות בכיוונים מנוגדים — עברית וערבית
מימין לשמאל (\foreignlanguage{english}{RTL}), לעומת אנגלית ורוב שפות אירופה
משמאל לימין (\foreignlanguage{english}{LTR}) \mbox{\cite{Unicode2023}}.

הבעיה מחריפה כאשר מסמך אקדמי מכיל מונחים טכניים באנגלית המוטמעים בתוך פסקאות
עבריות. מספרים, שמות של פונקציות תכנות ושמות ספריות — כולם מופיעים בתסריט
לטיני, ואילו הניתוח הסביבתי נכתב בעברית. ללא טיפול מדויק בעדיפות הכיווניות,
התוצאה עלולה להיות היפוך תווים, שבירת מילים ברצף הלא-נכון, והצגת מספרים
בסדר הפוך — כך שמספר כמו \foreignlanguage{english}{[10]} יוצג בטעות
כ-\foreignlanguage{english}{[01]}.

במערכות סוכנים מבוססות בינה מלאכותית, הדו-כיוניות מוסיפה שכבת מורכבות נוספת:
הפלט של סוכן אחד משמש קלט לסוכן הבא בשרשרת הייצור. אם הסוכן הכותב מייצר
טקסט עברי המכיל ביטויים טכניים באנגלית ללא סימון נכון, הסוכן העורך ואחר-כך
סוכן ייצור ה-\foreignlanguage{english}{LaTeX} לא יוכלו לנרמל את הכיוונים,
ובמסמך הסופי ייווצרו שגיאות נראות לעין שלא ניתן לתקנן בשלב מאוחר.
כדי להתמודד עם כך, על ממשק ה-\foreignlanguage{english}{SDK} לאכוף הנחיות
כיוון עקביות לאורך כל שלבי הצינור.
\end{hebrew}

% ---------------------------------------------------------------
\section[\texthebrew{פתרון \foreignlanguage{english}{LuaLaTeX} ו-\foreignlanguage{english}{Polyglossia}}]{\texthebrew{פתרון \foreignlanguage{english}{LuaLaTeX} ו-\foreignlanguage{english}{Polyglossia}}\mbox{\setLTR\enspace\thesection}}
\label{sec:bidi-solution}

\begin{hebrew}
מנוע \foreignlanguage{english}{LuaLaTeX} מספק תמיכה מלאה בתקן \foreignlanguage{english}{Unicode}
ברמת ה-\foreignlanguage{english}{engine}, כלומר ניתן לכתוב תווים עבריים ישירות בקבצי
המקור ללא קידוד מיוחד. בשילוב עם חבילת \foreignlanguage{english}{Polyglossia},
מקבלים פתרון מלא לניהול שפות מרובות במסמך אחד \mbox{\cite{Polyglossia2023}}.

ההגדרה הבסיסית דורשת הצהרת שפה ראשית ושפה משנית.
הפקודה \foreignlanguage{english}{\texttt{\textbackslash setmainlanguage\{hebrew\}}}
מגדירה עברית כשפת הבסיס של המסמך — כך שכל המסמך מוגדר כ-\foreignlanguage{english}{RTL}
כברירת-מחדל. הפקודה
\foreignlanguage{english}{\texttt{\textbackslash setotherlanguage\{english\}}}
מוסיפה אנגלית כשפה משנית; חלקים לטיניים מסומנים במפורש בפקודת
\foreignlanguage{english}{\texttt{\textbackslash foreignlanguage\{english\}\{...\}}}.

עבור מערכת הסוכנים המתוארת בעבודה זו, סוכן הכותב מייצר פסקאות עבריות שבתוכן
מונחים טכניים מסומנים בתגית ייעודית
\foreignlanguage{english}{\texttt{[EN:term]}}. שלב
עיבוד לאחר-מכן מחליף תגיות אלה בפקודת
\foreignlanguage{english}{\texttt{\textbackslash foreignlanguage\{english\}\{term\}}}
הנדרשת על-ידי \foreignlanguage{english}{LuaLaTeX}. גישה זו מאפשרת לסוכן הכותב
להתמקד בתוכן הסמנטי מבלי לדאוג לפרטי הסינטקס של \foreignlanguage{english}{LaTeX},
ומבטיחה כיווניות תקינה לאורך כל הצינור \mbox{\cite{Segal2026}}.
\end{hebrew}

% ---------------------------------------------------------------
\section[\texthebrew{אתגרי ייצור ופתרונות}]{\texthebrew{אתגרי ייצור ופתרונות}\mbox{\setLTR\enspace\thesection}}
\label{sec:bidi-challenges}

\begin{hebrew}
בפריסה של מערכות סוכנים בסביבת ייצור, האתגר הדו-כיווני מתבטא בארבעה תחומים עיקריים.

\textbf{ראשית}, שמות ספריות וסביבות-עבודה — כגון \foreignlanguage{english}{CrewAI},
\foreignlanguage{english}{LangChain} ו-\foreignlanguage{english}{LangGraph} — מופיעים
בתוך טקסט עברי ויוצרים ביטויים דו-כיווניים שמנוע
\foreignlanguage{english}{LuaLaTeX} חייב לפתור \mbox{\cite{CrewAI2024}}.
הפתרון הוא עטיפה עקבית של כל ביטוי לטיני בתוך
\foreignlanguage{english}{\texttt{\textbackslash foreignlanguage\{english\}\{...\}}},
בין אם מדובר בשמות מוצרים, ראשי-תיבות, כתובות ממשק או ערכי קוד.

\textbf{שנית}, ציטוטים ביבליוגרפיים שנוצרים בתוך סביבת
\foreignlanguage{english}{\texttt{\textbackslash begin\{hebrew\}}}
עלולים לגרום לשגיאת כיווניות כאשר מספר הציטוט הוא דו-ספרתי. מספר כמו
\foreignlanguage{english}{[12]} מוצג בטעות כ-\foreignlanguage{english}{[21]}.
הפתרון הוא עטיפה בפקודת
\foreignlanguage{english}{\texttt{\textbackslash mbox\{\textbackslash cite\{key\}\}}},
המבטיחה שהתיבה כולה מטופלת כיחידה \foreignlanguage{english}{LTR}
\mbox{\cite{Anthropic2024b}}.

\textbf{שלישית}, פקודות
\foreignlanguage{english}{\texttt{\textbackslash texttt\{...\}}}
בתוך בלוקים עבריים מייצרות היפוך של תווי \foreignlanguage{english}{ASCII},
משום ש-\foreignlanguage{english}{LuaLaTeX} מפרש את הארגומנטים בהקשר
ה-\foreignlanguage{english}{RTL} הפעיל.
הפתרון הוא עטיפה כפולה:
\foreignlanguage{english}{\texttt{\textbackslash foreignlanguage\{english\}\{\textbackslash texttt\{...\}\}}}.
כלל זה חל על כל פקודת \foreignlanguage{english}{LaTeX} המוצגת כקוד בתוך הטקסט העברי.

\textbf{רביעית}, כותרות פרקים וסעיפים בעברית מופיעות בתוכן-עניינים בסדר הפוך
אם לא מוגדר פורמט תוכן-עניינים מתאים. שימוש בחבילת
\foreignlanguage{english}{etoc} עם הגדרות \foreignlanguage{english}{RTL}
מותאמות-אישית — הכוללות נקודות-הנחיה עבריות (\foreignlanguage{english}{hebrewdotleaders})
וסדר הפוך של שם הסעיף ומספרו — פותר את הבעיה, כפי שניתן לראות בתוכן-העניינים
של עבודה זו.
\end{hebrew}

\addtocontents{toc}{\protect\etocstandardlines}\defaultheadingformat

% ---------------------------------------------------------------
\section{Mixed-Language Example}
\label{sec:bidi-mixed}

The following table catalogues the six most critical BiDi rules that production agent systems
must enforce when generating bilingual \LaTeX{} documents.
Each row states the problem category, the correct solution technique,
and a representative command example.
Agent prompts in this project encode these rules verbatim so that every writer-agent output
conforms before the compilation pipeline executes \cite{Segal2026}.

\begin{longtable}{@{}p{3.8cm}p{5.8cm}p{4.0cm}@{}}
\caption{BiDi Compliance Rules for Bilingual \LaTeX{} Generation}
\label{tab:bidi-rules}\\
\toprule
\textbf{Category} & \textbf{Solution Technique} & \textbf{Command Example} \\
\midrule
\endfirsthead
\toprule
\textbf{Category} & \textbf{Solution Technique} & \textbf{Command Example} \\
\midrule
\endhead
\midrule
\multicolumn{3}{r}{\footnotesize\itshape Continued on next page\ldots}\\
\endfoot
\bottomrule
\endlastfoot
English terms in Hebrew body &
Wrap every Latin word with
\texttt{\textbackslash foreignlanguage\{english\}\{Word\}} inside
\texttt{\textbackslash begin\{hebrew\}\ldots\textbackslash end\{hebrew\}} blocks;
applies to proper nouns, acronyms, library names, and API identifiers &
\texttt{\textbackslash foreignlanguage\{english\}\{LuaLaTeX\}} \\[4pt]
Citation digit reversal &
Wrap every \texttt{\textbackslash cite\{key\}} with
\texttt{\textbackslash mbox\{\ldots\}} so that multi-digit reference numbers
such as \texttt{[10]} are not reversed to \texttt{[01]} in RTL context &
\texttt{\textbackslash mbox\{\textbackslash cite\{key\}\}} \\[4pt]
Monospace code in Hebrew body &
Double-wrap: outer
\texttt{\textbackslash foreignlanguage\{english\}\{\ldots\}}
enforces LTR direction; inner
\texttt{\textbackslash texttt\{\ldots\}} applies monospace font;
omitting the outer wrapper mirrors every ASCII character in the PDF &
\texttt{\textbackslash foreignlanguage\{english\}\{\textbackslash texttt\{cmd\}\}} \\[4pt]
RTL section headings &
Enclose Hebrew title in \texttt{\textbackslash texthebrew\{\ldots\}};
append \texttt{\textbackslash mbox\{\textbackslash setLTR\textbackslash enspace\textbackslash thesection\}}
so that the section number renders in LTR order alongside the RTL heading &
\texttt{\textbackslash mbox\{\textbackslash setLTR\textbackslash enspace\textbackslash thesection\}} \\[4pt]
Table-of-contents RTL style &
Customise \texttt{etoc} chapter and section styles with Hebrew-aware
dot leaders (\texttt{\textbackslash hebrewdotleaders}) and reversed
\texttt{\textbackslash etocname}/\texttt{\textbackslash etocnumber} ordering;
restore standard lines after the Hebrew chapter with
\texttt{\textbackslash etocstandardlines} &
\texttt{\textbackslash hebrewdotleaders} \\[4pt]
Language declaration &
Declare \texttt{hebrew} as the primary document language so the entire
document defaults to RTL; register \texttt{english} as the other language
via \texttt{Polyglossia} so that LTR spans are available on demand &
\texttt{\textbackslash setmainlanguage\{hebrew\}} \\[4pt]
\end{longtable}

Failure to apply any single rule from Table~\ref{tab:bidi-rules} produces visible
typesetting defects — reversed digits, mirrored command names, or misaligned
section numbers — that are detectable by automated PDF inspection.
The rules were derived from the Unicode bidirectionality standard \cite{Unicode2023}
and the \texttt{Polyglossia} package documentation \cite{Polyglossia2023}.